% =============================================================================
% s7_experiments.tex --- §7 Experiments, COMPRESSED for the 9-page S version.
% Merges and compresses four long drafts:
%   drafts/s71_c0_results.tex  (C0 decision surface)
%   drafts/s7_c1_results.tex   (E3 / E7 / E10 / E8-E9-E1-E2-E12 ablations)
%   drafts/s7_c3_boundary.tex  (E5 melanoma net-negative + E6 severe + retake)
%   drafts/s77_dca_triage.tex  (DCA + global triage honest negative)
% Target: <= 1100 words / 2.0 pages. Sentences state key result + trend only;
%   numbers live in \input tables / figure captions / appendices.
% caveman OFF. All numbers copied verbatim from STORY locked table + source
%   drafts (verifier-checked); none invented or recomputed (red line 4).
%
% \input tables (paths relative to main_iclr9.tex / main_accv.tex in ICLR2027/):
%   table1_e10_sota.tex (tab:e10)        -- used in §7.4 (relocated to A28).
%   drafts_accv/tab_c0_surface.tex (tab:c0-surface) -- §7.1, replaces the
%       per-dimension reliability/recoverability prose.
%   drafts_accv/tab_ablations.tex (tab:ablations)   -- §7.6, replaces the
%       FiLM/cross-attn/PSNR/speed ablation prose.
% NOT \input here (intentional): table1_main.tex (tab:main, the demoted
%   calibration table -- footnoted per STORY R13, not part of the C0-C3
%   diagnosis-preservation experiments) and table2_universality.tex
%   (tab:universality, 5-backbone scaling -- on the firewall-1 universality
%   dimension, cite-not-reuse). Flagged for orchestrator.
%
% Cross-referenced labels DEFINED here (else upstream \cref undef):
%   sec:exp-surface, sec:exp-preserve, sec:exp-dploss, sec:exp-sota,
%   sec:exp-ablations, sec:exp-boundary, sec:exp-triage, sec:exp-dca.
% =============================================================================

\section{Experiments}\label{sec:exp}

We evaluate on the \itb{} dermoscopy benchmark, probing diagnostic content with a
frozen EfficientNet-B3 classifier kept \emph{independent} of the enhancer to avoid
circularity. Unless noted, every quantity carries a bootstrap $2000$-sample $95\%$
confidence interval, and significance means the CI excludes zero.

\subsection{The Reliability\,$\times$\,Recoverability Decision Surface}
\label{sec:exp-surface}
We chart the reliability\,$\times$\,recoverability surface introduced in
\cref{sec:decision-surface} over a $5\times5$ grid. Full per-cell AUC and recoverability
grids with CIs are reported in \cref{app:surface}.
Reliability degrades near-monotonically with severity. Blur is by far the most fragile
dimension and contrast the most robust. Recoverability splits into three regimes:
enhancement reliably recovers blur and colour shift from S3 onward, stays statistically
inconclusive on completeness, and significantly hurts exactly one of the $25$ cells,
contrast at the extreme level. Contrast thus looks robust yet is recoverable in name only,
so the gate must be dimension-aware. This single contrast-S5 cell is one per-dimension
trigger that complements the per-class evidence of \cref{sec:exp-boundary}. We do
\emph{not} generalise to ``enhancement is unsafe under severe degradation''.

\subsection{Diagnostic Preservation}\label{sec:exp-preserve}
On the moderate band the feature-level diagnosis-preserving model leaves diagnosis
essentially intact. The enhanced inputs reach a diagnostic $\Delta\mathrm{AUC}$ of
$-0.012$, within the $1.5\%$ target, with $95.8\%$ enhanced-versus-reference agreement and
a McNemar test finding no significant difference between enhanced and reference predictions
at $p=0.573$. The enhanced inputs are therefore diagnostically interchangeable with their
clean references on aggregate AUC. This claim is strictly confined to the
\emph{moderate-degradation window} measured on \emph{aggregate AUC}. The malignant
per-class boundary of \cref{sec:exp-boundary} is exactly what motivates the
query-for-retake gate. Aggregate preservation within the window and net-negative salvage on
the malignant subgroup are therefore not in tension but a two-regime division of labour:
enhance where the band holds, re-acquire where it does not.

\subsection{DP-Loss Ablation and the Mutual-Information Bound}\label{sec:exp-dploss}
Ablating the diagnosis-preserving objective \eqref{eq:dploss} against an otherwise
identical enhancer trained without it shows DP-Loss is what preserves diagnosis. On matched
bootstrap resamples it raises the diagnostic $\Delta\mathrm{AUC}_{\text{enh}}$ by $+0.0205$
with a CI excluding zero, lowers the diagnostic KL by $-0.148$ with a CI excluding zero, and
yields McNemar $p=2.3\times10^{-45}$. This is the empirical counterpart of the
mutual-information lower bound in Lemma~\ref{lem:dp}: tightening the feature-level KL budget
measurably raises the diagnostic information retained, the direction the bound predicts.
Together with the Stage-1 PSNR$\,\geq30$\,dB non-vacuity condition, this carries the
information-preservation claim of \cref{sec:prop3lemma3}. We do \emph{not} rely on a net
entropy-reduction claim, which the data do not support.

\subsection{Comparison with Generic Restoration SOTA}\label{sec:exp-sota}
Against six off-the-shelf non-diffusion restoration models applied zero-shot,
\visienhance{} is the only method that both restores fidelity and preserves diagnosis
(\cref{tab:e10}). Every generic model is a significantly worse diagnostic preserver, with
paired $\Delta\mathrm{AUC}$ CIs all excluding zero and McNemar $p<10^{-150}$. This
comparison is run under a \emph{matched zero-shot protocol} in which no method, ours
included, is adapted to the degraded distribution, so the gap is the
diagnostic-preservation gap between the objectives and not an artefact of one model having
seen the target degradations. A matched fine-tuned comparison (\todo{R1 核}) is reported in
\cref{sec:exp-transfer}, and the present claim is scoped to the zero-shot protocol.
Aggregate preservation does not, however, mean every lesion is safe. Enhancement still
pushes a small subset of correctly-called melanomas below the decision threshold, a
\emph{dangerous flip} shown in \cref{fig:enhancement}C, which we report transparently as
motivation for the query-for-retake channel (\cref{sec:agent}).

% E10 --- QP-Enhance vs 6 generic non-diffusion restoration SOTA (zero-shot).
% Frozen 会话28: job 1448952 (stored-mixed 降质, n=10881 paired, 117 melanoma; 与 E1 同
% degrade 算子 -> QP-Enhance PSNR 32.79 == E1 32.74 caliber). 数字 csv: results/e10_<m>.csv.
% 会话43 富表升级（BMVC 风格）: 彩色阴影 + per-column 粗体最优/下划线次优 + paired CI 列。
%   所有数字逐字搬自 STORY 锁定值 / ACCEPTANCE C1.4（results/e10_*.csv，verifier 核），零重算。
%   诚实红线（STORY L248 / ACCEPTANCE L183 dflip 单指标陷阱）：QP-Enhance dflip 0.194
%   非最优，绝不标粗/不标赢；dflip 列粗体给 Uformer 0.054（最优）、下划线给 Restormer 0.059。
% caveman OFF.
\begin{table}[!ht]
\centering
\caption{\textbf{E10 --- diagnosis preservation versus six generic restoration SOTA, zero-shot.}
All methods restore the identical mixed-severity degraded inputs, $n{=}10{,}881$ paired with $117$
melanomas. \textbf{Bold}/\underline{underline}: best/second per column; green: where \visienhance{}
wins. \visienhance{} wins five of six metrics, and every baseline is a significantly worse diagnostic
preserver: the paired block reports all $95\%$ CIs excluding $0$ with McNemar $p<10^{-150}$.
It does \emph{not} win dflip; actively rewriting within the salvage band occasionally flips a
melanoma, the motivation for query-for-retake (\cref{sec:agent}). The lowest-dflip baselines are also
the blurriest, with the highest KL and lowest consistency.}
\label{tab:e10}
\footnotesize
\setlength{\tabcolsep}{4.5pt}
\renewcommand{\arraystretch}{1.05}
\begin{tabular}{l l c c c c c c}
\toprule
Method & Pretraining & PSNR\,$\uparrow$ & SSIM\,$\uparrow$ & $\Delta$AUC\,$\to0$ & Consist.\,$\uparrow$ & KL\,$\downarrow$ & dflip\,$\downarrow$ \\
\midrule
\rowcolor{green!12}\textbf{\visienhance{} (FiLM+DP)} & dermoscopy
 & \cellcolor{green!18}\textbf{32.79}
 & \cellcolor{green!18}\textbf{0.907}
 & \cellcolor{green!18}\textbf{$-$0.017}
 & \cellcolor{green!18}\textbf{0.955}
 & \cellcolor{green!18}\textbf{0.10}
 & 0.194 \\
\midrule
Restormer \citep{zamir2022restormer}   & real denoising  & 22.30 & 0.836 & \cellcolor{red!16}$-$0.096 & 0.684 & 0.82 & \underline{0.059} \\
NAFNet \citep{chen2022simple}          & SIDD denoising  & 22.04 & 0.825 & \cellcolor{red!20}$-$0.115 & 0.669 & 0.88 & 0.077 \\
MIRNet-v2 \citep{zamir2022mirnetv2}    & LOL enhancement & 13.48 & 0.775 & \cellcolor{red!14}$-$0.087 & 0.818 & 0.50 & 0.158 \\
SwinIR \citep{liang2021swinir}         & colour denoising& 21.69 & 0.786 & \cellcolor{red!22}$-$0.134 & 0.821 & 0.36 & 0.478 \\
Uformer-B \citep{wang2022uformer}      & SIDD denoising  & 22.28 & \underline{0.836} & \cellcolor{red!16}$-$0.094 & 0.683 & 0.83 & \textbf{0.054} \\
Real-ESRGAN \citep{wang2021realesrgan} & $\times2$ SR    & 21.61 & 0.818 & \cellcolor{red!12}\underline{$-$0.083} & \underline{0.851} & \underline{0.35} & 0.234 \\
\midrule
\multicolumn{8}{l}{\textit{Paired $\Delta\mathrm{AUC}_{\text{baseline}-\text{VE}}$ vs.\ \visienhance{} (negative = baseline worse; $95\%$ CI; all exclude $0$):}}\\
\multicolumn{8}{l}{\;\; Restormer $-0.079\,[-0.096,-0.064]$\quad NAFNet $-0.098\,[-0.116,-0.080]$\quad MIRNet-v2 $-0.070\,[-0.088,-0.054]$}\\
\multicolumn{8}{l}{\;\; SwinIR $-0.116\,[-0.135,-0.098]$\quad Uformer-B $-0.077\,[-0.093,-0.061]$\quad Real-ESRGAN $-0.066\,[-0.082,-0.050]$}\\
\bottomrule
\end{tabular}
\end{table}
% tab:e10 --- SOTA-comparison centrepiece, relocated back into §7.4 body (会话43).

\begin{figure}[t]
  \centering
  \includegraphics[width=0.74\linewidth]{figures/fig4_boundary_agent.pdf}
  \caption{\textbf{The honest boundary and the query-for-retake gate.}
  \textbf{(A)} Per-class salvage on distinct conditional pools: enhancement recovers most
  degradation-misclassified benign lesions ($75.6\%$, $1809/2392$) but only $5.2\%$ of the
  salvageable melanoma pool ($4/77$) while damaging $31.0\%$ of the at-risk melanoma pool
  ($85/274$), a net loss of malignant cases left unchanged by mask-weighted retraining; see
  \cref{sec:exp-boundary} for the full per-class denominators. \textbf{(B)} Decision-curve and triage analysis on
  \itb{}-LQ ($n=300$): net benefit is statistically indistinguishable across direct diagnosis and
  all variants (overlapping $95\%$ CIs), and at a $20\%$ referral budget direct diagnosis attains
  the highest auto-handled sensitivity ($0.818$) while the strongest variant only tracks it
  ($0.788$). \textbf{(C)} The retake fraction rises monotonically with degradation
  ($0.055$/$0.651$/$0.889$ across the high, moderate, and severe bands).}
  \label{fig:boundary-agent}
  \label{fig:dca}% alias: appendix \cref{fig:dca} resolves here (DCA panel B) in the short version
\end{figure}

% Fig 3 (enhancement vs restoration) relocated to Appendix A28 (app:moved-figs) for the 9-page
% limit; fig:enhancement label lives there, main-text \cref{fig:enhancement} resolves to the appendix.

\subsection{DP-Loss is Transferable}\label{sec:exp-transfer}
To separate the objective from our particular network, we graft DP-Loss onto a strong
generic restoration backbone with its architecture unchanged, and fine-tune each generic
baseline on degraded dermoscopy so the comparison no longer credits \visienhance{} for
adaptation alone (\cref{sec:exp-sota}). The grafting test (\todo{R3 核 $\Delta$AUC}, $95\%$
CI \todo) decides whether the diagnosis-preserving effect is attributable to the objective
rather than the architecture; the per-class melanoma picture (\todo{R3 核}, \todo) is
reported in full and not aggregated away, and should the grafted backbone remain
net-negative on melanoma we attribute the residual failure to the restoration boundary
itself, consistent with \cref{sec:exp-boundary}. All quantities here are pending the
fine-tuning and grafting runs, and the direction of the effect --- and hence the verbs ---
is fixed only once they complete (\todo{R1/R3 核 — HPC 排队中，未填即占位}).

\subsection{Conditioning, Fidelity, and Speed Ablations}\label{sec:exp-ablations}
Quality-conditional FiLM is a diagnostic, not a perceptual, mechanism. Removing it leaves
fidelity essentially unchanged yet degrades diagnostic preservation on every axis, and a
heavier cross-attention module is statistically indistinguishable from it, so we keep FiLM
on parsimony grounds (\cref{tab:ablations}). A per-degradation breakdown in
\cref{fig:enhancement}B confirms the enhancer is strongest on brightness and blur but
genuinely \emph{harmful} on contrast, consistent with the contrast-S5 trigger of
\cref{sec:exp-surface}.
% =============================================================================
% tab_ablations.tex --- §7.6 conditioning, fidelity, and speed ablations.
% Compresses the number-dense §7.6 prose into one ablation table.
% caveman OFF. All numbers copied verbatim from STORY_FRAMEWORK / ACCEPTANCE
%   locked rows E1/E8/E9/E12 and the existing §7.6 prose; none recomputed
%   (red line 4). Sources: results/e1_film_ablation.json (E1),
%   results/filmabl_diag.json (E8), results/stage2_diag_paired_e9.csv (E9),
%   results/e12_speed.csv (E12). Single-seed point estimates; paired-bootstrap
%   2000-sample 95% CIs and McNemar p where stated.
% \input from main_accv.tex / main_iclr9.tex (path relative to that file).
% =============================================================================
\begin{table}[!ht]
\centering
\caption{\textbf{Conditioning, fidelity, and speed ablations.}
Quality-conditional FiLM is a diagnostic, not a perceptual, mechanism: removing it leaves
fidelity essentially unchanged yet degrades diagnostic preservation on every axis. A heavier
cross-attention module is statistically indistinguishable from FiLM, so we keep FiLM on
parsimony grounds. PSNR $32.74$\,dB clears the $30$\,dB non-vacuity gate of
Lemma~\ref{lem:dp}. Best diagnostic-preservation value in bold.}
\label{tab:ablations}
\vspace{4pt}
\footnotesize
\setlength{\tabcolsep}{5pt}
\renewcommand{\arraystretch}{1.05}
\begin{tabular}{lccccc}
\toprule
Configuration & $\Delta$AUC\,$\to0$ & Agreement\,$\uparrow$ & Diag.\ KL\,$\downarrow$ & PSNR (dB)\,$\uparrow$ & Params \\
\midrule
\rowcolor{green!12}With FiLM (ours) & \cellcolor{green!18}$\mathbf{-0.033}$ & \cellcolor{green!18}$\mathbf{0.90}$ & \cellcolor{green!18}$\mathbf{0.24}$ & $32.74$ & --- \\
Without FiLM            & $-0.042$          & $0.87$          & $0.35$          & $\mathbf{33.06}$ & --- \\
Cross-attention ($+1.8$M) & \multicolumn{3}{c}{indistinguishable (CIs incl.\ 0, McNemar $p{=}0.679$)} & --- & $+1.8$M \\
\bottomrule
\end{tabular}
\vspace{2pt}
\footnotesize
\emph{Deterministic enhancer (per image):} PSNR $32.74$\,dB, SSIM $0.91$ ($n{=}19878$);
runtime $16.08$\,ms/image (p95 $17.0$).
\vspace{-4pt}
\end{table}


\subsection{The Honest Boundary: When Enhancement Hurts Diagnosis}
\label{sec:exp-boundary}
A per-class split (\cref{fig:boundary-agent}A, with full denominators in its caption) is
where the salvage story collapses. The benign and malignant subgroups behave in opposite
directions: enhancement salvages most degradation-misclassified benign lesions while
yielding a \emph{net loss of correctly-called melanomas}. We report the two melanoma rates
on distinct conditional pools by construction. Salvage is measured over the
\emph{salvageable pool} already misclassified after degradation, and damage over the
\emph{at-risk pool} still correctly called, so the two denominators differ because the
events are conditioned on opposite pre-enhancement outcomes and not because cases were
selected. Because the benign subgroup dominates the population, the benign-dominated
aggregate moderate-band salvage rate is \emph{not} a meaningful capability number, and we
deliberately do not report it as one. Only the per-class breakdown reflects clinical risk.
This is not a reconstruction-loss artefact. Retraining with a mask-weighted objective
leaves melanoma salvage an honest null (\cref{fig:boundary-agent}A), so the failure mode
lies beyond what per-pixel $L_1$ can reach. The boundary is sharper still on the severe
band, where enhancement significantly \emph{lowers} diagnostic AUC by
$\Delta\mathrm{AUC}=-0.0559$ with a $95\%$ CI of $[-0.085,-0.028]$. The agent accordingly
routes severe and high-risk inputs to query-for-retake rather than enhance, with the retake
fraction rising strictly with degradation across bands (\cref{fig:boundary-agent}C),
realising the threshold-window policy of Theorem~\ref{thm:agent-compact}. These per-class,
per-severity, and per-dimension triggers are the hardest motivation for the gate, but they
are \emph{local, per-band} statements only. The honest global comparison is in
\cref{sec:exp-triage}.

\subsection{Cost--Benefit, Decision Curves, and the Global Triage Negative Result}
\label{sec:exp-triage}
\label{sec:exp-dca}% upstream \cref{sec:exp-dca} (rem:no-global) resolves here
On \itb{}-LQ, the worst-case low-quality subset with $n=300$, we run net-benefit
decision-curve analysis and a confidence-gated triage simulation
(\cref{fig:boundary-agent}B; cost model in \cref{app:a20}). The gate needs only an
off-the-shelf per-input uncertainty signal, and a standard variational information
bottleneck suffices.\footnote{The quality-conditional variant we also report carries
\emph{no} method-gain claim. It appears only as the strongest tested configuration, and it
still does not outperform direct diagnosis.} \textbf{Net benefit does not separate the
methods.} Across the clinically relevant threshold range the maximum net benefit is
statistically indistinguishable among direct diagnosis and all bottleneck variants, spanning
$0.179$ to $0.192$ with all $95\%$ CIs overlapping, and the only positive statement is that
each stays above the treat-all reference. \textbf{The global triage comparison is an honest
negative.} At a $20\%$ referral budget direct diagnosis attains the \emph{highest}
auto-handled sensitivity of $0.818$, the strongest bottleneck variant only tracks it at
$0.788$ without exceeding it, and under-confident variants collapse. Global quality-aware
triage therefore does not outperform direct diagnosis, and the agent's global risk is not
guaranteed to be at most that of direct diagnosis. This is the honest global counterpart of
the band-local guarantee: Theorem~\ref{thm:agent-compact} bounds risk only \emph{within} the
moderate window, as stated in Remark~\ref{rem:no-global}, and the result delimits where the
system should defer to re-acquisition rather than assert a diagnosis. The defensible positive
reading is narrow. A calibrated posterior supports reliable confidence-driven abstention, and
closing the band-local-to-global gap by \emph{learning} the routing thresholds is left to
future work.
% (DCA/triage figure merged into Fig.~\ref{fig:boundary-agent}B; retake gradient into panel C.)
